\documentclass[10pt,a4paper]{amsart}
\usepackage[margin=19mm]{geometry}
\usepackage{amsmath,amssymb,mathtools}
\usepackage[scr=rsfs,cal=euler]{mathalfa}
\usepackage{xfrac}
\usepackage[unicode,hidelinks,bookmarks=false]{hyperref}
\newcommand{\ie}{i.e.\ }
\newcommand{\cf}{cf.\ }
\newcommand{\resp}{respectively\ }
\newcommand{\Subsection}{\subsection}
\providecommand{\nobreakdash}{\nobreak\hbox{-}}
\setlength{\emergencystretch}{2em}
\multlinegap0pt
\usepackage{bm}
\usepackage{bbm}
\usepackage{dsfont}
\usepackage{xspace}
\usepackage{cancel}
\usepackage{mathtools}
\usepackage{amsthm}
\makeatletter
\@ifundefined{theorem}{\newtheorem{theorem}{Theorem}}{}
\@ifundefined{corollary}{\newtheorem{corollary}[theorem]{Corollary}}{}
\@ifundefined{definition}{\newtheorem{definition}[theorem]{Definition}}{}
\@ifundefined{lemma}{\newtheorem{lemma}[theorem]{Lemma}}{}
\@ifundefined{proposition}{\newtheorem{proposition}[theorem]{Proposition}}{}
\makeatother
\def\C{\Gamma}
\def\L{\Lambda}
\def\NN{{\mathbb N}}
\def\RR{{\mathbb R}}
\def\ae{\text{-a.e.}\ }
\def\a{\alpha}
\def\c{\gamma}   \def\C{\Gamma}
\def\d{\delta}   %\def\C{\Gamma}
\def\e{\epsilon}
\def\pX{{\partial X}}
\def\p{{\partial}}
\newcommand{\CCC}{\mathcal{C}}
\newcommand{\GGG}{\mathcal{G}}
\title[Patterson-Sullivan construction of equilibrium states and we]{Patterson-Sullivan construction of equilibrium states and weighted counting in nonpositive curvature}
\author{Weisheng Wu}
\date{Source: arXiv:2509.09924v2 (CC BY 4.0)}
\begin{document}
\maketitle

\noindent\emph{Source and licence.} Patterson-Sullivan construction of equilibrium states and weighted counting in nonpositive curvature. Version arXiv:2509.09924v2 (CC BY 4.0), distributed under the Creative Commons Attribution 4.0 International licence, as recorded in the arXiv licence metadata for this exact version and shown on the abstract page https://arxiv.org/abs/2509.09924v2. Full rights notice: licenses/arxiv-2509.09924-CC-BY-4.0.txt.\par\smallskip

\noindent\textit{Editorial scope.} Counted unit: the Proposition whose source label is \texttt{Busemann2} in the article (source0.tex lines 1070--1074, source label \texttt{Busemann2}), delivered together with the article's own complete original proof of it (source0.tex lines 1075--1131, one \texttt{proof} environment of 57 lines). No mathematics is written, repaired or re-derived: statement and proof are the article's own text line for line. Required context retained uncounted: smalldelta (lines 417--419); corelimit (lines 470--474); limitzero (lines 507--509); decay (lines 534--538); Bowen1 (lines 753--757); cocexi2 (lines 766--768); full (lines 1013--1015). Every definition, display and label of those lines is retained unchanged. Omitted from this package and not redistributed: 1--416, 420--469, 475--506, 510--533, 539--752, 758--765, 769--1012, 1016--1069, 1132--2623. Nothing inside a retained excerpt is omitted or altered: every retained body is delivered exactly as the article's lines are written, so no omission receipt of any kind accompanies this package. Citations: the retained excerpts carry no citation command at all, so no bibliography entry of the article is transcribed and no reference is redistributed. Reference adaptation: none was needed and none was made; every reference inside the delivered unit resolves inside it, so no unresolved reference or citation is produced. Printed slips kept literal: no printed word, symbol, hypothesis or number of the counted statement or of its proof was altered. Layout and class: the article's own class is \texttt{amsart} with packages including amssymb, amsmath, paralist, graphics, epsfig, graphicx, epstopdf, enumitem, comment, hyperref; the wrapper is the lane's pinned \texttt{amsart} editorial wrapper at 10pt on A4, which restates the theorem-like environments the retained text needs and adds the article's own macro definitions that the retained text needs (\texttt{\textbackslash C}, \texttt{\textbackslash CCC}, \texttt{\textbackslash GGG}, \texttt{\textbackslash L}, \texttt{\textbackslash NN}, \texttt{\textbackslash RR}, \texttt{\textbackslash a}, \texttt{\textbackslash ae}, \texttt{\textbackslash c}, \texttt{\textbackslash d}, \texttt{\textbackslash e}, \texttt{\textbackslash p}, \texttt{\textbackslash pX}), with extra wrapper packages (bm, bbm, dsfont, xspace, cancel, mathtools). The delivered numbering is the unit's own. Assets: no retained span contains a figure environment, a graphics-inclusion command, a PDF-page inclusion, a table or a TikZ picture; no image file is copied into this package, the package declares no asset dependency and no image file is redistributed with it. Licence: version arXiv:2509.09924v2 of this article is distributed under the Creative Commons Attribution 4.0 International licence, as recorded in the arXiv licence metadata for this exact version and shown on the abstract page https://arxiv.org/abs/2509.09924v2. A full rights notice is delivered at licenses/arxiv-2509.09924-CC-BY-4.0.txt. Characters: every retained excerpt is pure ASCII, so the delivered engine, XeLaTeX with its default Latin Modern text font, reports no missing character.
\begin{equation}\label{smalldelta}
d_K(v,w)< \delta e^\Lambda \Rightarrow |\lambda(v)-\lambda(w)|\le \eta/3.    
\end{equation}

\begin{definition}\label{corelimit}
Let $\CCC$ denote the set of geodesic segments $[o,\c o]$ from $o$ to $\c o$ where $\c\in \C$. Given $L>0$, let $\L_c^L(\C)$ be the set of $\xi\in \p X$ such that there exists $\c_n\in \C$ such that $[o,\c_no]\in \GGG^L$ and $\xi=\lim_{n\to \infty}\c_no.$
Then we define the \emph{core limit set} as
$$\L_c(\C):=\bigcup_{L>0}\L_c^L(\C).$$
\end{definition}

\begin{lemma}\label{limitzero}
Let $v=\dot{c}_{o,\xi}(0)$ where $\xi\in \L_c^L(\C)$. If $p\in H^s(v)$, then $\lim_{t\to +\infty}d^s(\dot c_{p,\xi}(t), \dot c_{o,\xi}(t))=0$.
\end{lemma}

\begin{lemma}\label{decay}
Let $v=\dot{c}_{o,\xi}(0)$ where $\xi\in \L_c^L(\C)$. If $p\in X$, then 
$$d_K(\dot c_{p,\xi}(t+T_3), \dot c_{o,\xi}(t))\le \d e^{\frac{\eta}{2}T_4}e^{-\frac{\eta}{2}t}, \ \forall t>T_5,$$
for some $T_3(p,\xi)\in \RR$ and $0<T_4(p,\xi) <T_5(p,\xi)$ independent of $t>0$.
\end{lemma}

\begin{lemma}\label{Bowen1}
Let $v=\dot{c}_{o,\xi}(0)$ where $\xi\in \L_c^L(\C)$. For any $p\in X$ and $T>T_5$, we have
$$\left|\int_p^{c_{p,\xi}(T+T_3)}F-\int_{o}^{\pi g^Tv}F\right|\le K\d^\a e^{\frac{\eta\a}{2}T_4}\int_{T_5}^{T}e^{-\frac{\eta\a}{2}t}dt+\|F\|(d(p,o)+2T_5)$$
where $T_3(p,\xi)$, $T_4(p,\xi)$ and $T_5(p,\xi)$ are from Lemma \ref{decay} and independent of $T$.
\end{lemma}

\begin{corollary}\label{cocexi2}
If $p, q\in X$ and $\xi\in \L_c^L(\C)$, then Gibbs cocycle $C_{F,\xi}(q,p)$ is well defined.    
\end{corollary}

\begin{proposition}\label{full}
    $\mu_{F,o}((\L_c(\C))^c)=0$.
\end{proposition}

\begin{proposition}\label{Busemann2}
For every $\c\in \C$, we have
$$\frac{d\mu_{F,\c o}}{d\mu_{F,o}}(\xi)=e^{-C_{F-\d_F, \xi}(\c o,o)}$$
for $\mu_{F,o} \ae\ \xi\in \pX$.
\end{proposition}

\begin{proof}
By Proposition \ref{full}, it is sufficient to prove the proposition for $\xi\in \L_c$. So the Busemann cocycle $C_{F-s, \xi}(\c o,o)$ is well defined by Corollary \ref{cocexi2} for any $s\ge \d_F$.

Fix $\d$ from \eqref{smalldelta}. By Lemma \ref{limitzero} and Corollary \ref{cocexi2}, for any small $\e>0$, there exists $T>0$ large enough such that 
$$d^s(\dot c_{o,\xi}(t), \dot c_{\c o,\xi}(t+s_0))<\d/10, \quad \forall t\ge T$$ 
and
\begin{equation}\label{e:est1}
\begin{aligned}
\left|C_{F-s,\xi}(\c o,o)-\int_{o}^{c_{o,\xi}(T)}(F-s)+\int_{\c o}^{c_{\c o,\xi}(T+s_0)}(F-s)\right|<\e/5
\end{aligned}
\end{equation}
where $s_0=b_\xi(\c o,o)$. Since $\xi\in \L_c$, let $\c_n\in \C$ be as in Definition \ref{corelimit}. Then $\c_n o\to \xi$ as $n\to \infty$. For small $0<\rho \ll \min\{\e, \d/10\}$, pick $n$ large enough such that 
$$\max\{\angle_o(\c_n o, c_{o,\xi}(T)), \angle_{\c o}(\c_n o, c_{\c o,\xi}(T+s_0))\}\ll \rho.$$
Therefore by comparison theorem,
$$\max\{d_K(\dot c_{o,\c_n o}(T), \dot c_{o,\xi}(T)), d_K(\dot c_{\c o,\c_n o}(T+s_0), \dot c_{\c o,\xi}(T+s_0))\}< \rho.$$
Therefore, if $\rho$ is small enough, then
\begin{equation}\label{e:est2}
\begin{aligned}
\left|\int_{o}^{c_{o,\xi}(T)}(F-s)-\int_{o}^{c_{o,\c_n o}(T)}(F-s)\right|&<\e/5\\
\left|\int_{\c o}^{c_{\c o,\xi}(T+s_0)}(F-s)-\int_{\c o}^{c_{o,\c_n o}(T+s_0)}(F-s)\right|&<\e/5.
\end{aligned}
\end{equation}

By the choice of $T$ and $\rho$, $d_K(\dot c_{o,\c_n o}(T),\dot c_{\c o,\c_n o}(T+s_0))<\d/5$.
Recall that $[o,\c o]\in \GGG^L$. By a similar argument as in the proof of Lemma \ref{decay}, there exists $T_1>T$ such that for any $T_1<t<d(o,\c_no)-L$,
\[\frac{1}{t-T}\int_T^t\lambda(g^s \dot c_{o,\c_n o}(0))ds\ge \frac{99\eta}{100}.\]
Then following a similar argument as in the proof of Lemma \ref{Bowen1} and noting that $L$ is fixed, we have
\begin{equation}\label{e:est3}
\begin{aligned}
\left|\int_{c_{o,\c_n o}(T_1)}^{\c_n o}(F-s)-\int_{c_{\c o,\c_n o}(T_1+s_0)}^{\c_n o}(F-s)\right|<\e/5
\end{aligned}
\end{equation}
by enlarging $T$ if necessary. But we still need consider the time interval $[T,T_1]$. Here we emphasize that the choice of $T_1$ is universal for all orbit segments in $\GGG^L$, and hence independent of the particular orbit segment $[o,\c_n o]$. Thus we can enlarge $n$ enough, so that  \eqref{e:est1}, \eqref{e:est2} still hold for $T_1$ instead of $T$, and meanwhile on this new orbit segment $[o,\c_n o]$, \eqref{e:est3} also holds.
Combining \eqref{e:est1}, \eqref{e:est2} (both with $T_1$ instead of $T$) and \eqref{e:est3}, we know
\begin{equation*}
\begin{aligned}
C_{F-s,\xi}(\c o,o)=\lim_{n\to \infty }\int_{o}^{\c_n o}(F-s)-\int_{\c o}^{\c_n o}(F-s).
\end{aligned}
\end{equation*}

Now take a sequence of neighborhoods $\{U_i\}_{i=1}^\infty$ of $\xi$ in $\overline X$ such that $\{\xi\}=\cap_{i} U_i$ and $\mu_{F,o}(\partial U_i)=\mu_{F,\c o}(\partial U_i)=0, \forall i\in \NN$. We have showed that for any $\c_i o\in U_i$,
\begin{equation*}
\begin{aligned}
\frac{e^{\int_{\c o}^{\c_i o}(F-s)}}{e^{\int_{o}^{\c_i o}(F-s)}}\to e^{-C_{F-s,\xi}(\c o,o)}
\end{aligned}
\end{equation*}
as $i\to \infty$. From the above proof, this convergence is uniform for all $s$ in a compact interval $[\d_F, \d_F+c]$ for some $c>0$.
%Note that $\lim_{r\to +\infty}\frac{h(t+r)}{h(r)}=1$ for any $t>0$. 
Then
\begin{equation*}
\begin{aligned}
&\frac{d\mu_{F,\c o}}{d\mu_{F,o}}(\xi)=\lim_{i\to \infty}\frac{\mu_{F, \c o}(U_i)}{\mu_{F, o}(U_i)}=\lim_{i\to \infty}\lim_{s_k\searrow \d_F}\frac{\mu_{F, \c o,s_k}(U_i)}{\mu_{F, o,s_k}(U_i)}\\
=&\lim_{s_k\searrow \d_F}\lim_{i\to \infty}\frac{\mu_{F, \c o,s_k}(U_i)}{\mu_{F, o,s_k}(U_i)}=e^{-C_{F-\d_F,\xi}(\c o, o)}.
\end{aligned}
\end{equation*}
The proof of the proposition is complete.
\end{proof}

\end{document}
